\section{Discussion and Future Work}

We have built, through the addition of one generator and three axioms,
an elegant and complete diagrammatic theory for Parametric Linear Programming.
We argue that such a result not only unlocks a new way to interact with LP problems in an algebraic
and computational way, but also provides a new perspective on optimization through the lens of
a compositional framework. This lifts optimization, and LP problems in particular, from a functional setting to a relational-theoretic framework within the Extended-Lawvere Quantale category.

The development followed a strict linear discipline, adding exactly one new
generator at a time and verifying at each step that the resulting algebra is
sound, expressive, and complete with respect to a precisely defined semantic
category. The fact that such a linear development is
possible --- that each new class of optimization problems requires only a single
additional generator and a finite set of axioms --- is evidence that the
compositional structure is not imposed artificially but present in the
underlying mathematics. Polyhedral relations and quadratic relations are not merely two
classes of optimization problems that happen to admit graphical representations;
they are parallel enrichments of the same underlying affine structure, each adding
a different kind of expressive power, and their common base is a theorem rather than
an accident of notation.

The $\mathbf{GLP}$ theory has a natural extension to an even more expressive theory
when enriched with more powerful generators. We believe that the class
of quadratic programs could be, just as was done in this work, fully axiomatized by such a
broader SMT. Another direction is the exploration of the functoriality of the bifunction
adjoint (the Legendre--Fenchel transform extended to morphisms)~\cite{willerton2015legendre, touchette2005legendre}: this suggests a duality theory
analogous to LP duality should be expressible diagrammatically as well.

A more applied direction concerns the computational implications of diagrammatic
normal forms. If two optimization problems are equal as morphisms in
$\mathbf{LPRel}$ if and only if their diagrams are related by the rewriting rules
of $\mathbf{GLP}$, then diagram simplification is problem simplification: a
rewriting system on string diagrams is simultaneously a symbolic preprocessor
for optimization solvers. Such an application of graphical algebras to theorem
proving was already pointed to in other works \cite{bonchi2021.9farkas,bonchi2021polyhedral}.
This work is a step toward fully expressing optimization within
its own algebraic setting.

